\scanpage{42}
% archive foliation: 41
\[
\delta=\rho^N\delta\qquad\forall N\in\mathbb{N}
\]
d'où, prenant $N$ grand d'où $\rho^N=0$ :
\[
\delta=0\qquad\text{cqfd.}
\]
\uline{Extension aux tenseurs}
\[
\begin{aligned}
u&\in\operatorname{End}_K(M)\\
v&\in\operatorname{End}_K(N)
\end{aligned}
\]
Alors $M$ $(N)$ $\overline G_a$-Module par $\bar e_u$ $(\bar e_v)$, et on a $M\otimes N$ $\overline G_a$-Module par $\bar e_u\otimes\bar e_v$ --- et
\[
\bar e_u\otimes\bar e_v=\bar e_{u\otimes\operatorname{id}_N+\operatorname{id}_M\otimes v}
\]
De $=$, si $u,v$ nilpotents
\[
e_u\otimes e_v=e_{u\otimes\operatorname{id}_N+\operatorname{id}_M\otimes v}
\]
Extension : $\displaystyle\bigotimes_{i\in I}M_i$ ($I$ ens d'indices fini, de cardinal $\geq2$)

\uline{Corollaire} Soit $M$ muni d'une structure d'alg. (pas néc. assoc. ni unitaire, encore moins comm.) sur $K$, $u\in\operatorname{End}_{K\text{-mod}}(M)$, donc $\overline G_a$ (Resp. $G_a$ si $u$ nilpotent) opère sur le $K$-module $M$. Pour que cette opération soit une op. d'endom d'algèbre, il f et s. que $u$ soit une \uline{dérivation}.
